\documentclass[10pt, a4paper, twoside]{article}
% =========================================================
% Packages
% =========================================================

\usepackage[utf8]{inputenc}

% Mathematics
\usepackage{amsmath,amssymb,amsthm}

% Springer-like fonts
% \usepackage{newtxtext}
% \usepackage{newtxmath}

% Page layout
\usepackage[
  paperwidth=160mm,
  paperheight=240mm,
  inner=20mm,
  outer=16mm,
  top=18mm,
  bottom=21mm,
  headheight=14pt,
  headsep=6mm,
  footskip=10mm
]{geometry}

% Headers / footers
\usepackage{fancyhdr}

% Graphics / diagrams
\usepackage{graphicx}
\usepackage{tikz}

\usetikzlibrary{
  positioning,
  arrows.meta
}

% Lists / columns
\usepackage{enumitem}
\usepackage{multicol}
\usepackage{paracol}

% Colors / framed environments
\usepackage{xcolor}
\usepackage[framemethod=tikz]{mdframed}

% Paragraph formatting
\usepackage{indentfirst}

% Dummy text for testing
\usepackage{blindtext}

% Index
\usepackage{imakeidx}

% Hyperref should be loaded late
\usepackage[hidelinks]{hyperref}

% Code listings
\usepackage{minted}

\setminted{
    bgcolor=gray!8,
    frame=single,
    framesep=4mm,
    baselinestretch=1.05,
    fontsize=\small,
    breaklines=true
}


% =========================================================
% Page / Typography Setup
% =========================================================

\definecolor{sepiabg}{RGB}{244,236,220}

% Uncomment if desired:
% \pagecolor{sepiabg}

% Paragraphs
\setlength{\parindent}{1.5em}
\setlength{\parskip}{0pt}

% Slightly tighter Springer-like display mathematics
\setlength{\abovedisplayskip}{7pt plus 2pt minus 2pt}
\setlength{\belowdisplayskip}{7pt plus 2pt minus 2pt}
\setlength{\abovedisplayshortskip}{4pt plus 2pt}
\setlength{\belowdisplayshortskip}{5pt plus 2pt minus 2pt}

% Slightly compact lists
\setlist{
  topsep=4pt,
  itemsep=1.5pt,
  parsep=0pt,
  partopsep=0pt
}


% =========================================================
% Headers / Footers
% =========================================================

\pagestyle{fancy}
\fancyhf{}

% Even pages: current section
\fancyhead[LE]{%
  \small\itshape
  \nouppercase{\leftmark}
}

% Odd pages: book / course title
\fancyhead[RO]{%
  \small\itshape
  Honor Algebra
}

% Centered page number
\fancyfoot[C]{%
  \small\thepage
}

\renewcommand{\headrulewidth}{0pt}
\renewcommand{\footrulewidth}{0pt}

% Section name in running header
\renewcommand{\sectionmark}[1]{%
  \markboth{\thesection\enspace #1}{}
}


% =========================================================
% Section Behaviour
% =========================================================

% Every section begins on a new page
\let\oldsection\section

\renewcommand{\section}{%
  \clearpage
  \oldsection
}


% =========================================================
% Equation Numbering
% =========================================================

\numberwithin{equation}{subsection}


% =========================================================
% Index
% =========================================================

\makeindex[
  intoc,
  title=Index
]


% =========================================================
% Shortcuts for Brackets
% =========================================================

\newcommand{\paren}[1]{\left(#1\right)}
\newcommand{\abs}[1]{\left|#1\right|}
\newcommand{\norm}[1]{\left\lVert#1\right\rVert}
\newcommand{\bracks}[1]{\left[#1\right]}
\newcommand{\set}[1]{\left\{#1\right\}}


% =========================================================
% Theorem Environments
% ONE numbered system
% =========================================================

\theoremstyle{plain}

\newtheorem{theorem}{Theorem}[section]
\newtheorem{lemma}[theorem]{Lemma}
\newtheorem{proposition}[theorem]{Proposition}
\newtheorem{corollary}[theorem]{Corollary}

\theoremstyle{definition}

\newtheorem{definition}[theorem]{Definition}
\newtheorem{example}[theorem]{Example}

\theoremstyle{remark}

\newtheorem*{remark}{Remark}


% =========================================================
% Cambridge-style Semantic Environments
% =========================================================

\theoremstyle{definition}

\newtheorem*{axiom}{Axiom}
\newtheorem*{assumption}{Assumption}
\newtheorem*{notation}{Notation}


% =========================================================
% Index Term Command
% =========================================================

\newcommand{\term}[1]{%
  \emph{#1}\index{#1}%
}


% =========================================================
% Lists
% =========================================================

\renewcommand{\labelitemi}{--}
\renewcommand{\labelitemii}{$\circ$}
\renewcommand{\labelenumi}{(\roman{enumi})}


% =========================================================
% QED Symbol
% =========================================================

\renewcommand{\qedsymbol}{%
  \ensuremath{\blacksquare}%
}

% Alternative:
% \renewcommand{\qedsymbol}{$\square$}


% =========================================================
% Custom Mathematical Commands
% =========================================================

\newcommand{\R}{\mathbb{R}}
\newcommand{\N}{\mathbb{N}}
\newcommand{\Z}{\mathbb{Z}}
\newcommand{\Q}{\mathbb{Q}}
\newcommand{\C}{\mathbb{C}}

\newcommand{\ds}{\displaystyle}

% =========================================================
% Color Commands
% =========================================================

\newcommand{\blue}[1]{%
  \textcolor{blue}{#1}%
}

\newcommand{\red}[1]{%
  \textcolor{red}{#1}%
}


% =========================================================
% Figure Helper
% =========================================================

\newcommand{\mypic}[3]{%
  \begin{figure}[htbp]
    \centering
    \includegraphics[width=#3\textwidth]{#1}
    \caption{#2}
  \end{figure}
}


% =========================================================
% Word Definition Environment
% =========================================================

\newenvironment{worddef}[1]
  {%
    \par
    \noindent
    \textbf{#1:}\ 
    \itshape
  }
  {%
    \par
    \normalfont
  }


% =========================================================
% Title
% =========================================================

\title{%
  \textbf{Math 135}
  \\[0.5em]
  \large Honor Algebra
}

\author{Qinghao Hu}

\date{\today}


% =========================================================
% Compilation Control
% =========================================================

% Uncomment and edit to only compile a subset of chapters while drafting:
% \includeonly{chapters/c1}


% =========================================================
% Document
% =========================================================

\begin{document}

\maketitle

\clearpage

\tableofcontents

\section{Admin}
%!TEX root = ../../note

\subsection{General Information}

\subsubsection{Contact}
For mathematics questions: \texttt{abeltaos@uwaterloo.ca}

For course logistics: \texttt{math135@uwaterloo.ca}

\subsubsection{Teaching Style}
The LEARN website contains several useful resources. In particular:
\begin{enumerate}
    \item The icons at the top of the page provide quick access to important resources, including the course schedule.
    \item Crowdmark is used for submitting assignments.
\end{enumerate}

\subsubsection{Graded Assignments}
There are three assignments each week.

\subsubsection{Late Submission Policy}
Late submissions are accepted with a penalty of $1\%$ per minute for the entire assignment.

\subsubsection{Office Hours}
Monday and Wednesday 2:00 to 3:30 in (MC 6503) or by appointment (email the instructor to arrange). 

%!TEX root = ../../note

\section{Introduction to the Language of Mathematics}
%!TEX root = ../../note

\subsection{Introduction to Proof}
Checking a few examples may suggest that a statement is true. A proof
explains why it holds for every object covered by the statement.

\begin{definition}[Proof]
A proof is a logically valid argument that establishes the truth of a mathematical statement using definitions, axioms, and previously proved results.
\end{definition}

\begin{proposition}
For every positive integer $n$, the number $n^2 + 1$ is not a perfect square.
\end{proposition}

\begin{proof}
Let $n$ be a positive integer. The consecutive perfect squares surrounding $n^2 + 1$ are $n^2$ and $(n+1)^2$. Indeed,
\begin{equation*}
    n^2 < n^2 + 1 < n^2 + 2n + 1 = (n+1)^2,
\end{equation*}
where the second inequality follows from $n > 0$.

Therefore, $n^2 + 1$ lies strictly between two consecutive perfect squares, so it cannot itself be a perfect square.
\end{proof}

\subsubsection*{In practice}
\begin{itemize}
    \item \textbf{Ruling out a perfect square.} Locate the number strictly
    between two consecutive perfect squares. Check both inequalities and
    identify where the hypothesis on the number is used.
\end{itemize}

%!TEX root = ../../../note
\subsection{Sets}
\begin{definition}[Set]
    A \textbf{set} is a well-defined, unordered collection of distinct objects. Each object that appears 
    in this collection is called an \textbf{element} of the set. 
\end{definition}

For a finite set, we can denote it by listing its elements, separated by commas,
surrounding the entire list by braces (or by $\set{}$).  

\begin{example}
    The set containing the elements 2, 7, and -3 can be denoted by 
    \begin{equation*}
        \set{7, -3, 2}. 
    \end{equation*}
\end{example}

\begin{example}
    Here are some examples:
    \begin{itemize}
        \item $\set{2, 4, 6, 8}$ 
        \item $\set{1, 2, \set{1, 2, 3}}$
        \item $\set{a, b, c, d}$
        \item $\N = \set{1, 2, 3, \cdots}$
        \item $\Z = \set{\cdots, -1, 0, 1, \cdots}$
        \item $\Q = \set{\frac{a }{b } : a, b \in \Z, b \neq 0}$
    \end{itemize}
    Fact: $\abs{\N} = \aleph_0, \quad \abs{\R} = \aleph_1 \quad \aleph_1 = 2^{\aleph_0}$. 
\end{example}

\begin{definition}
    The \textbf{parity} of an integer is \textbf{even} if the integer is a multiple of 2, and \textbf{odd} otherwise. 
\end{definition}

\subsection{Statement and Negation}
\begin{definition}[Statement]
    A \textbf{statement} is a sentence that has a definite state of being either true or false. 
\end{definition}

\begin{definition}[Negation]
    Suppose that $A$ is a statement. Then the negation of $A$, denoted by $\neg A$, is the statement 
    asserting the opposite truth value to $A$. That is, $\neg A$ is false when $A$ is true, and $\neg A$ is true 
    when $A$ is false. 
\end{definition}


% Uncomment when you want the index printed:
% \printindex

\end{document}